(1) 理想弱连接的直流电流—相位关系为 \(I_s=I_c\sin\phi\)，相位演化满足 \(\hbar\,d\phi/dt=2eV\)，其中电荷 \(2e\) 来自Cooper对。前式描述零电压下的超电流，后式说明非零电压会使两侧相对相位连续随时间改变。
(2) 固定电压 \(V\) 时，\(\phi(t)=\phi_0+2eVt/\hbar\)，所以超电流 \(I_s(t)=I_c\sin(\phi_0+\omega_Jt)\)，\(\omega_J=2eV/\hbar\)，普通频率 \(f_J=\omega_J/(2\pi)=2eV/h\)。频率与弱连接的具体几何无关，是电压与量子常数的直接联系。
(3) 令相位自由能满足 \(I_s=(2e/\hbar)\,dU_J/d\phi\)，代入 \(I_c\sin\phi\) 得 \(dU_J/d\phi=(\hbar I_c/(2e))\sin\phi\)。积分得到 \(U_J(\phi)=-E_J\cos\phi+\mathrm{const}\)，其中 \(E_J=\hbar I_c/(2e)\)，单位为能量，最小值在相位差为 \(2\pi\) 整数倍时取得。
(4) 对称两结超导环在忽略环自感时，相位差满足 \(\phi_1-\phi_2=2\pi\Phi/\Phi_0\)，\(\Phi_0=h/(2e)\)。总电流 \(I=I_c(\sin\phi_1+\sin\phi_2)=2I_c\sin[(\phi_1+\phi_2)/2]\cos(\pi\Phi/\Phi_0)\)，优化平均相位得 \(I_{c,\rm SQUID}(\Phi)=2I_c|\cos(\pi\Phi/\Phi_0)|\)，对磁通具有 \(\Phi_0\) 周期。
